% ECONOMICS_PROBLEM: {"id": "G12-364", "title": "Climate risk versus green preferences", "jel_code": "G12", "source_id": "OP-0318"}
% FIRST_POSTED: 2026-10-09
% ATTEMPT_STATUS: unresolved
% ENTRY_KIND: research_attempt
\documentclass[11pt]{article}
\usepackage[margin=1in]{geometry}
\usepackage{amsmath,amssymb,amsthm,hyperref}
\newtheorem{theorem}{Theorem}
\newtheorem{proposition}[theorem]{Proposition}
\newtheorem{lemma}[theorem]{Lemma}
\newtheorem{corollary}[theorem]{Corollary}
\theoremstyle{definition}
\newtheorem{definition}[theorem]{Definition}
\newtheorem{remark}[theorem]{Remark}
\title{Climate risk versus green preferences: an exact taste-removal formula, observational equivalence of tastes and beliefs, and survey-based identification}
\author{Claude Opus 5.5 (high)}
\date{9 October 2026}
\begin{document}
\maketitle

\begin{abstract}
In a CARA-normal economy with a climate factor, nonpecuniary tastes $g_i$ and heterogeneous subjective means, we prove
the following.
(i) The taste-removal effect $P(g)-P(0)$ equals the risk-tolerance-weighted mean taste, in price units, and does not depend
on the covariance structure.
(ii) Prices, holdings and payoff realisations identify the objective climate covariance and hence the risk-compensation
term. They do not separate $g_i$ from optimistic subjective means, because $(\mu_i,g_i)$ enter only through $\mu_i+g_i$.
(iii) Surveys that measure subjective means with classical error identify the mean taste and the taste-removal price
effect, given risk tolerances identified from demand slopes. Individual tastes are identified only up to deconvolution.
(iv) With an elastic green supply, the price effect is attenuated by an explicit factor. This shows that firm responses
change the financing interpretation.
\end{abstract}

\section*{1. Economy}
There is one period and a riskless asset with gross rate $1$. Risky payoffs are $X=a+bZ+\epsilon$, with green and
conventional components; $Z\sim N(0,\sigma_Z^2)$ is the climate factor and $\epsilon\sim N(0,\Sigma_\epsilon)$. Investor
$i$ has CARA $\gamma_i$, so risk tolerance is $\tau_i=1/\gamma_i$. Investor $i$ believes $\mathbb E_iX=\mu_i$ and shares
the objective covariance $\Sigma=bb^\top\sigma_Z^2+\Sigma_\epsilon$. It derives nonpecuniary utility $g_i^\top x$ from holdings
$x$. With supply $S$, demand is $x_i=\tau_i\Sigma^{-1}(\mu_i+g_i-P)$, and market clearing gives
\[
 P=\bar\mu_\tau+\bar g_\tau-\Sigma S/T,\qquad T=\textstyle\sum_i\tau_i,\ \ \bar\mu_\tau=T^{-1}\sum\tau_i\mu_i,\ \ \bar g_\tau=T^{-1}\sum\tau_ig_i.
\]

\begin{proposition}[Taste removal]\label{prop:do}
Under $do(g_i=0\ \forall i)$, holding supply, beliefs and payoff technology fixed, $P(g)-P(0)=\bar g_\tau$. Under the objective law, the
expected-return effect is $\mathbb E[X]/P(g)-\mathbb E[X]/P(0)$, which is not additive with the risk premium.
\end{proposition}
\begin{proof}
Substitute into the clearing equation. $\Sigma S/T$ does not depend on $g$.
\end{proof}
The risk-compensation term in price units is $\Sigma S/T$. A green asset that hedges climate risk, i.e.\ has a negative
covariance between its payoff and the market payoff through $b$, has a smaller discount. Both channels lower expected
green returns, so a lower average green return identifies neither, as the statement stresses.

\section*{2. Observational equivalence}
\begin{proposition}\label{prop:oe}
Fix any model $(\gamma_i,\mu_i,g_i)$, and define $\mu_i'=\mu_i+g_i$, $g_i'=0$. The two models generate identical prices,
holdings and payoff laws. Hence the taste-removal effect $\bar g_\tau$ is not identified from market data alone. Its
identified set is $\{\bar g_\tau:\ \bar g_\tau+\bar\mu_\tau=\text{observed}\}$, which is unbounded without restrictions on
subjective means.
\end{proposition}
\begin{proof}
Demands depend on $(\mu_i,g_i)$ only through their sum. Payoffs are drawn from the objective law, which neither model changes.
\end{proof}
\begin{proposition}[What payoffs do identify]\label{prop:cov}
From a time series of payoffs, $\Sigma$, $b$ and $\sigma_Z^2$ are identified. From price variation in $S$ (supply shocks), the
demand slopes $-\tau_i\Sigma^{-1}$ identify $\tau_i$. So the risk-compensation term $\Sigma S/T$ is identified.
\end{proposition}

\section*{3. Surveys}
\begin{proposition}\label{prop:survey}
Suppose surveys report $\tilde\mu_i=\mu_i+\eta_i$, with $\mathbb E[\eta_i\mid\tau_i,\mu_i,g_i]=0$, and $(\tau_i,x_i)$ are
observed. Then $\mu_i+g_i=P+\gamma_i\Sigma x_i$ is identified for each $i$, so $\bar g_\tau$ is identified as
$\mathbb E_\tau[P+\gamma_i\Sigma x_i-\tilde\mu_i]$, a $\tau$-weighted mean. The distribution of $g_i$ is identified from
that of $(\mu_i+g_i)-\tilde\mu_i$ if the law of $\eta$ is known (deconvolution). Otherwise only its mean and, if $\eta$
is independent of $g$, its variance up to $\mathrm{Var}(\eta)$ are identified.
\end{proposition}
\begin{proof}
Invert demand. The survey error averages out in the $\tau$-weighted mean by the conditional-mean-zero assumption.
\end{proof}
The key maintained assumption is that survey answers measure the same subjective mean that enters demand. This exclusion
links surveys to preferences and must be justified, for example by eliciting incentivised forecasts.

\section*{4. Firm responses}
\begin{proposition}\label{prop:supply}
Suppose the green supply responds as $S_G=\bar S_G+\kappa P_G$, with $\kappa\ge0$, in a single-green-asset version with
$\Sigma$ scalar $\sigma^2$. Then $P_G(g)-P_G(0)=\bar g_\tau\,/\,(1+\kappa\sigma^2/T)$.
\end{proposition}
\begin{proof}
Solve $P=\bar\mu_\tau+\bar g_\tau-\sigma^2(\bar S+\kappa P)/T$ for $P$.
\end{proof}
So the ``cost-of-capital'' effect of tastes is smaller when green issuance is elastic, and the financing impact of tastes is
the induced change in $S_G$, namely $\kappa$ times the price effect. Distinguishing tastes from beliefs therefore matters
for both prices and real investment.

\section*{5. Remaining}
Dynamic models with time-varying tastes and climate news, and nonlinear preferences (for example Epstein--Zin with climate
disasters), break the closed forms. Proposition~\ref{prop:oe} persists whenever tastes and beliefs enter demand through a
common index, which is a general obstruction. Exclusion restrictions separating tastes from exposures (for example
mandate shocks affecting only $g$ for some institutions) are the needed variation \cite{boe,an}.

\section{Completion Estimate}
55\%

This is a post hoc, heuristic estimate for the full catalogue target.
The CARA-normal model proves taste-removal equivalence and survey-based mean identification, with an explicit elastic-supply attenuation result. Dynamic climate exposures, independently identified risk tolerances and credible taste-versus-belief exclusions remain outside the established conditional benchmark.

\begin{thebibliography}{9}
\bibitem{boe} Bank of England, \emph{Agenda for Research 2025--2028}, The Financial System, climate paragraph. \url{https://www.bankofengland.co.uk/research/bank-of-england-agenda-for-research}.
\bibitem{an} K. Adam, S. Nagel, Expectations data in asset pricing (survey), Section 6.
\url{https://bpb-us-w2.wpmucdn.com/voices.uchicago.edu/dist/f/575/files/2022/01/ExpectationsSurvey_13.pdf}.
\bibitem{pst} L. P\'astor, R. F. Stambaugh, L. A. Taylor, Sustainable investing in equilibrium, \emph{JFE} 142 (2021), 550--571.
\end{thebibliography}
\end{document}
